% ============================================================
%  ШАБЛОН ОТЧЁТА
%  Лабораторная работа №1
%  Дисциплина: Информационно-поисковые системы
%  РУДН, 2025
% ============================================================

\documentclass[12pt, a4paper]{report}

% ── Кодировки и язык ────────────────────────────────────────
\usepackage[utf8]{inputenc}
\usepackage[T2A]{fontenc}
\usepackage[russian, english]{babel}

% ── Геометрия страницы (ГОСТ 7.32) ──────────────────────────
\usepackage[
  left   = 30mm,
  right  = 15mm,
  top    = 20mm,
  bottom = 20mm
]{geometry}

% ── Шрифты ──────────────────────────────────────────────────
\usepackage{mathptmx}          % Times New Roman (основной текст)
\usepackage{helvet}            % Helvetica (заголовки)
\usepackage[scaled=0.85]{beramono} % моноширинный шрифт

% ── Межстрочный интервал и отступы ──────────────────────────
\usepackage{setspace}
\setstretch{1.5}               % полуторный интервал (ГОСТ)

\usepackage{indentfirst}
\setlength{\parindent}{1.25cm}
\setlength{\parskip}{0pt}

% ── Заголовки разделов ───────────────────────────────────────
\usepackage{titlesec}

\titleformat{\chapter}[block]
  {\sffamily\Large\bfseries\color{NavyBlue}}
  {\thechapter.}{0.5em}{}
\titlespacing{\chapter}{0pt}{12pt}{8pt}

\titleformat{\section}[block]
  {\sffamily\large\bfseries\color{SteelBlue}}
  {\thesection.}{0.5em}{}
\titlespacing{\section}{0pt}{10pt}{6pt}

\titleformat{\subsection}[block]
  {\sffamily\normalsize\bfseries\color{SteelBlue!80!black}}
  {\thesubsection.}{0.5em}{}
\titlespacing{\subsection}{0pt}{8pt}{4pt}

% ── Цвета ────────────────────────────────────────────────────
\usepackage[dvipsnames, table]{xcolor}
\definecolor{NavyBlue}   {RGB}{31,  78, 121}
\definecolor{SteelBlue}  {RGB}{46, 117, 182}
\definecolor{LightBlue}  {RGB}{222,234,246}
\definecolor{PlaceholderBg}{RGB}{255,251,230}
\definecolor{PlaceholderBorder}{RGB}{200,180,120}
\definecolor{CodeBg}     {RGB}{245,245,245}
\definecolor{CodeBorder} {RGB}{46, 117, 182}

% ── Таблицы ─────────────────────────────────────────────────
\usepackage{booktabs}
\usepackage{tabularx}
\usepackage{multirow}
\usepackage{longtable}
\usepackage{array}
\renewcommand{\arraystretch}{1.4}

% ── Листинги кода ────────────────────────────────────────────
\usepackage{listings}
\lstset{
  language        = Python,
  basicstyle      = \ttfamily\small,
  keywordstyle    = \bfseries\color{NavyBlue},
  commentstyle    = \itshape\color{gray},
  stringstyle     = \color{SteelBlue!80!black},
  backgroundcolor = \color{CodeBg},
  frame           = leftline,
  framerule       = 2pt,
  rulecolor       = \color{CodeBorder},
  numbers         = left,
  numberstyle     = \tiny\color{gray},
  numbersep       = 8pt,
  showstringspaces= false,
  breaklines      = true,
  breakatwhitespace= true,
  inputencoding   = utf8,
  extendedchars   = true,
  literate        = % поддержка кириллицы в lstlisting
    {а}{{\cyra}}1 {б}{{\cyrb}}1 {в}{{\cyrv}}1 {г}{{\cyrg}}1
    {д}{{\cyrd}}1 {е}{{\cyre}}1 {ж}{{\cyrzh}}1 {з}{{\cyrz}}1
    {и}{{\cyri}}1 {й}{{\cyrishrt}}1 {к}{{\cyrk}}1 {л}{{\cyrl}}1
    {м}{{\cyrm}}1 {н}{{\cyrn}}1 {о}{{\cyro}}1 {п}{{\cyrp}}1
    {р}{{\cyrr}}1 {с}{{\cyrs}}1 {т}{{\cyrt}}1 {у}{{\cyru}}1
    {ф}{{\cyrf}}1 {х}{{\cyrh}}1 {ц}{{\cyrc}}1 {ч}{{\cyrch}}1
    {ш}{{\cyrsh}}1 {щ}{{\cyrshch}}1 {ъ}{{\cyrhrdsn}}1 {ы}{{\cyrery}}1
    {ь}{{\cyrsftsn}}1 {э}{{\cyrerev}}1 {ю}{{\cyryu}}1 {я}{{\cyrya}}1
    {А}{{\CYRA}}1 {Б}{{\CYRB}}1 {В}{{\CYRV}}1 {Г}{{\CYRG}}1
    {Д}{{\CYRD}}1 {Е}{{\CYRE}}1 {Ж}{{\CYRZH}}1 {З}{{\CYRZ}}1
    {И}{{\CYRI}}1 {К}{{\CYRK}}1 {Л}{{\CYRL}}1 {М}{{\CYRM}}1
    {Н}{{\CYRN}}1 {О}{{\CYRO}}1 {П}{{\CYRP}}1 {Р}{{\CYRR}}1
    {С}{{\CYRS}}1 {Т}{{\CYRT}}1 {У}{{\CYRU}}1 {Ф}{{\CYRF}}1
    {Х}{{\CYRH}}1 {Ц}{{\CYRC}}1 {Ч}{{\CYRCH}}1 {Ш}{{\CYRSH}}1
    {Щ}{{\CYRSHCH}}1 {Ъ}{{\CYRHRDSN}}1 {Ы}{{\CYRERY}}1 {Ь}{{\CYRSFTSN}}1
    {Э}{{\CYREREV}}1 {Ю}{{\CYRYU}}1 {Я}{{\CYRYA}}1,
  captionpos      = b,
  xleftmargin     = 1.25cm,
}

% ── Рамки и боксы ────────────────────────────────────────────
\usepackage{tcolorbox}
\tcbuselibrary{skins, breakable}

% Рамка-заполнитель (placeholder)
\tcbset{
  placeholderbox/.style = {
    colback    = PlaceholderBg,
    colframe   = PlaceholderBorder,
    fonttitle  = \sffamily\small\bfseries,
    fontupper  = \itshape\small\color{gray!60!black},
    boxrule    = 0.6pt,
    left       = 6pt, right = 6pt,
    top        = 4pt, bottom= 4pt,
    breakable,
    before upper = {},
  },
  infobox/.style = {
    colback  = LightBlue,
    colframe = NavyBlue,
    fonttitle= \sffamily\small\bfseries\color{NavyBlue},
    boxrule  = 1pt,
    left     = 8pt, right = 8pt,
    top      = 6pt, bottom= 6pt,
    breakable,
  },
}

\newcommand{\placeholder}[2][Вставьте результат]{%
  \begin{tcolorbox}[placeholderbox, title={#2}]
    #1
  \end{tcolorbox}%
  \vspace{4pt}%
}

% ── Прочие пакеты ────────────────────────────────────────────
\usepackage{graphicx}
\usepackage{subcaption}
\usepackage{caption}
\captionsetup{font=small, labelfont={bf,sf}, labelsep=period}

\usepackage{enumitem}
\setlist[enumerate]{leftmargin=1.5cm, itemsep=2pt, parsep=0pt}
\setlist[itemize]  {leftmargin=1.5cm, itemsep=2pt, parsep=0pt}

\usepackage{hyperref}
\hypersetup{
  colorlinks = true,
  linkcolor  = NavyBlue,
  urlcolor   = SteelBlue,
  citecolor  = NavyBlue,
  pdftitle   = {Отчёт ЛР №1 — Информационно-поисковые системы},
}

\usepackage{fancyhdr}
\usepackage{lastpage}
\usepackage{amsmath}
\usepackage{microtype}

% ── Колонтитулы ──────────────────────────────────────────────
\pagestyle{fancy}
\fancyhf{}
\fancyhead[R]{\sffamily\footnotesize\color{gray}Лабораторная работа №1\ |\ Информационно-поисковые системы}
\fancyfoot[C]{\sffamily\footnotesize\color{gray}\thepage\ из\ \pageref{LastPage}}
\renewcommand{\headrulewidth}{0.4pt}
\renewcommand{\footrulewidth}{0.4pt}
\renewcommand{\headrule}{\color{SteelBlue}\hrule width\headwidth height\headrulewidth}
\renewcommand{\footrule}{\color{SteelBlue}\hrule width\headwidth height\footrulewidth}

% ── Подавление колонтитулов на chapter-страницах ─────────────
\fancypagestyle{plain}{%
  \fancyhf{}
  \fancyfoot[C]{\sffamily\footnotesize\color{gray}\thepage\ из\ \pageref{LastPage}}
  \renewcommand{\headrulewidth}{0pt}
  \renewcommand{\footrulewidth}{0.4pt}
  \renewcommand{\footrule}{\color{SteelBlue}\hrule width\headwidth height\footrulewidth}
}

% ── Нумерация глав без слова «Глава» ─────────────────────────
\renewcommand{\chaptername}{}

% ============================================================
\begin{document}
\setcounter{chapter}{0}

% ── Титульная страница ────────────────────────────────────────
\begin{titlepage}
\thispagestyle{empty}
\begin{center}

  {\sffamily\small\bfseries МИНИСТЕРСТВО НАУКИ И ВЫСШЕГО ОБРАЗОВАНИЯ\\
  РОССИЙСКОЙ ФЕДЕРАЦИИ}

  \vspace{6pt}
  {\sffamily\small\bfseries ФЕДЕРАЛЬНОЕ ГОСУДАРСТВЕННОЕ АВТОНОМНОЕ\\
  ОБРАЗОВАТЕЛЬНОЕ УЧРЕЖДЕНИЕ ВЫСШЕГО ОБРАЗОВАНИЯ}

  \vspace{8pt}
  {\sffamily\normalsize\bfseries
  «РОССИЙСКИЙ УНИВЕРСИТЕТ ДРУЖБЫ НАРОДОВ\\
  ИМЕНИ ПАТРИСА ЛУМУМБЫ»}

  \vspace{6pt}
  {\sffamily\small Факультет \underline{\hspace{6cm}}}\\[2pt]
  {\sffamily\small Кафедра \underline{\hspace{7cm}}}

  \vspace{10pt}
  \textcolor{SteelBlue}{\rule{\linewidth}{1pt}}
  \vspace{12pt}

  {\sffamily\Large\bfseries\color{NavyBlue}
  ИНФОРМАЦИОННО-ПОИСКОВЫЕ СИСТЕМЫ}

  \vspace{20pt}

  {\sffamily\LARGE\bfseries ОТЧЁТ}\\[6pt]
  {\sffamily\large по лабораторной работе №1}

  \vspace{12pt}
  \begin{tcolorbox}[infobox, title={}]
    \centering\sffamily\large\bfseries\color{NavyBlue}
    Предобработка текстовых данных\\
    для анализа судебных документов
  \end{tcolorbox}

  \vspace{30pt}

\end{center}

% Таблица данных студента
\begin{flushright}
  \begin{tabular}{p{5cm}p{7cm}}
    \toprule
    \rowcolor{LightBlue}
    \sffamily\small\bfseries Поле & \sffamily\small\bfseries Значение \\
    \midrule
    \sffamily\small Студент       & \underline{\hspace{6.5cm}} \\[4pt]
    \sffamily\small Группа        & \underline{\hspace{6.5cm}} \\[4pt]
    \sffamily\small Специальность & Информационно-управляющие системы \\[4pt]
    \sffamily\small Преподаватель & \underline{\hspace{6.5cm}} \\[4pt]
    \sffamily\small Дата сдачи    & \underline{\hspace{6.5cm}} \\[4pt]
    \sffamily\small Оценка        & \underline{\hspace{6.5cm}} \\
    \bottomrule
  \end{tabular}
\end{flushright}

\vfill
\begin{center}
  {\sffamily\normalsize Москва, 2025}
\end{center}
\end{titlepage}

% ── Оглавление ───────────────────────────────────────────────
\tableofcontents
\thispagestyle{plain}
\clearpage

% ============================================================
%  РАЗДЕЛ 1. ЦЕЛЬ И ЗАДАЧИ
% ============================================================
\chapter{Цель и задачи работы}

Целью настоящей лабораторной работы является освоение методов предобработки
текстовых данных применительно к корпусу судебных документов, а также формирование
практических навыков применения инструментов обработки естественного языка (NLP)
на языке программирования Python.

В ходе выполнения работы студент решает следующие задачи:

\begin{enumerate}
  \item Формирование и первичный анализ учебного корпуса судебных актов.
  \item Реализация конвейера очистки текстовых данных.
  \item Токенизация, нормализация регистра и фильтрация стоп-слов.
  \item Морфологическая нормализация (лемматизация) с применением библиотеки \texttt{pymorphy2}.
  \item Извлечение именованных сущностей (NER) с применением библиотеки \texttt{natasha}.
  \item Реализация ключевое-словарного классификатора документов.
  \item Количественный анализ корпуса и визуализация результатов.
\end{enumerate}

% ============================================================
%  РАЗДЕЛ 2. ТЕОРЕТИЧЕСКАЯ СПРАВКА
% ============================================================
\chapter{Краткая теоретическая справка}

\begin{tcolorbox}[infobox, title={Источники}]
  Подробное изложение теоретического материала представлено в конспектах Лекции 1
  («Введение в информационно-аналитические системы в юридической практике»)
  и Лекции 2 («Сбор и предобработка данных для анализа судебной практики»).
\end{tcolorbox}

\vspace{8pt}

\section{Предобработка текстовых данных}

Предобработка данных представляет собой совокупность методов, приводящих исходный
неструктурированный текст в форму, пригодную для применения алгоритмов машинного обучения.
Качество предобработки напрямую определяет качество последующей аналитики: неудалённые
артефакты форматирования, нелемматизированные словоформы или сохранённые стоп-слова снижают
точность моделей классификации~\cite{lec2}.

\section{Ключевые операции NLP-конвейера}

\textbf{Токенизация} — процесс разбиения текста на минимальные смысловые единицы (токены).
Является обязательным первым шагом любого NLP-конвейера.

\textbf{Лемматизация} — приведение слова к его нормальной словарной форме с учётом
морфологических свойств. Для русскоязычных текстов применяется библиотека \texttt{pymorphy2},
обеспечивающая морфологический анализ на основе словарей OpenCorpora.

\textbf{Стоп-слова} — частотные служебные слова (предлоги, союзы, местоимения),
не несущие самостоятельного смысла и удаляемые при анализе текста.

\textbf{Распознавание именованных сущностей} (NER, Named Entity Recognition) —
задача выделения из текста упоминаний объектов реального мира (персон, организаций,
дат) и их классификации по типам~\cite{lec1}. В данной работе используется библиотека
\texttt{natasha}, специализированная для обработки русскоязычных текстов.

% ============================================================
%  РАЗДЕЛ 3. ХОД ВЫПОЛНЕНИЯ
% ============================================================
\chapter{Ход выполнения работы}

\section{Анализ корпуса (Задание 1.1)}

Приведите результаты первичного анализа учебного корпуса: количество документов по
категориям, средний объём текста.

\placeholder[Вставьте таблицу или вывод программы с результатами анализа корпуса]{Результаты анализа корпуса}

% ── 3.2 ─────────────────────────────────────────────────────
\section{Очистка текста (Задания 2.1–2.2)}

Опишите реализованные функции очистки. Приведите исходный код и поясните каждый шаг.

\placeholder[Вставьте реализацию функций \texttt{normalize\_whitespace} и \texttt{remove\_punctuation} с комментариями]{Код функций очистки}

\placeholder[Вставьте пример: фрагмент текста до и после очистки]{Пример результата очистки}

% ── 3.3 ─────────────────────────────────────────────────────
\section{Токенизация и удаление стоп-слов (Задания 3.1–3.2)}

Обоснуйте выбор дополнительных юридических стоп-слов. Приведите результаты токенизации.

\placeholder[Перечислите добавленные юридические стоп-слова и обоснуйте выбор каждого]{Дополнительные юридические стоп-слова}

\placeholder[Вставьте пример вывода: первые 15 токенов после фильтрации для выбранного документа]{Пример токенизации}

% ── 3.4 ─────────────────────────────────────────────────────
\section{Лемматизация (Задания 4.1–4.2)}

Приведите реализацию функции \texttt{lemmatize\_tokens}. Продемонстрируйте, что
лемматизация объединяет словоформы одного слова.

\placeholder[Вставьте код функции \texttt{lemmatize\_tokens}]{Реализация функции лемматизации}

Для наглядности сравните частотные словари до и после лемматизации. Пример оформления
представлен в таблице~\ref{tab:lemma_compare}. Заполните её собственными данными.

\begin{table}[h]
  \centering
  \caption{Сравнение частотных словарей до и после лемматизации (топ-10)}
  \label{tab:lemma_compare}
  \rowcolors{2}{LightBlue!40}{white}
  \begin{tabularx}{0.9\textwidth}{c X c X c}
    \toprule
    \rowcolor{LightBlue}
    \textsf{\textbf{\#}} &
    \textsf{\textbf{Токен}} &
    \textsf{\textbf{Частота}} &
    \textsf{\textbf{Лемма}} &
    \textsf{\textbf{Частота}} \\
    \midrule
    1  & \textit{[слово]} & — & \textit{[лемма]} & — \\
    2  & \textit{[слово]} & — & \textit{[лемма]} & — \\
    3  & \textit{[слово]} & — & \textit{[лемма]} & — \\
    4  & \textit{[слово]} & — & \textit{[лемма]} & — \\
    5  & \textit{[слово]} & — & \textit{[лемма]} & — \\
    6  & \textit{[слово]} & — & \textit{[лемма]} & — \\
    7  & \textit{[слово]} & — & \textit{[лемма]} & — \\
    8  & \textit{[слово]} & — & \textit{[лемма]} & — \\
    9  & \textit{[слово]} & — & \textit{[лемма]} & — \\
    10 & \textit{[слово]} & — & \textit{[лемма]} & — \\
    \bottomrule
  \end{tabularx}
\end{table}

% ── 3.5 ─────────────────────────────────────────────────────
\section{Извлечение именованных сущностей (Задания 5.1–5.2)}

Приведите результаты работы NER-модуля. Укажите наблюдения относительно точности.

\placeholder[Вставьте вывод программы: именованные сущности для документа по выбору студента]{Именованные сущности в выбранном документе}

\placeholder[Вставьте таблицу: топ-5 организаций и персон по всему корпусу]{Сводная таблица именованных сущностей}

% ── 3.6 ─────────────────────────────────────────────────────
\section{Ключевое-словарная классификация (Задания 6.1–6.2)}

Опишите принцип работы реализованного классификатора. Укажите, какие леммы были
добавлены в словарь и на каком основании.

\placeholder[Вставьте финальную версию словаря ключевых лемм по категориям]{Расширенный словарь ключевых лемм}

Результаты классификации представлены в таблице~\ref{tab:classification}.

\begin{table}[h]
  \centering
  \caption{Результаты ключевое-словарной классификации документов}
  \label{tab:classification}
  \rowcolors{2}{LightBlue!30}{white}
  \begin{tabularx}{\textwidth}{l X X c}
    \toprule
    \rowcolor{LightBlue}
    \textsf{\textbf{ID документа}} &
    \textsf{\textbf{Истинная категория}} &
    \textsf{\textbf{Предсказанная категория}} &
    \textsf{\textbf{Верно}} \\
    \midrule
    A40-001-2023 & & & \\
    A40-002-2023 & & & \\
    A56-003-2023 & & & \\
    A56-004-2023 & & & \\
    A60-005-2023 & & & \\
    A60-006-2023 & & & \\
    A07-007-2023 & & & \\
    A07-008-2023 & & & \\
    A40-009-2023 & & & \\
    A40-010-2023 & & & \\
    \midrule
    \multicolumn{3}{r}{\textsf{\textbf{Точность}}} & \textbf{\underline{\hspace{1.5cm}}} \\
    \bottomrule
  \end{tabularx}
\end{table}

% ── 3.7 ─────────────────────────────────────────────────────
\section{Визуализация результатов (Задания 7.1–7.3)}

Приведите построенные графики и облака слов. Каждый рисунок сопровождается
подписью и пояснительным текстом.

\placeholder[Вставьте изображение диаграммы (команда \texttt{\textbackslash includegraphics})]{Рис.~1. Диаграмма распределения документов по категориям}

\placeholder[Вставьте облако слов по всему корпусу (лемматизированные токены)]{Рис.~2. Облако слов по всему корпусу}

\placeholder[Задание 7.3 (повышенная сложность): вставьте облака слов по отдельным категориям]{Рис.~3–6. Облака слов по категориям}

% ============================================================
%  РАЗДЕЛ 4. КОНТРОЛЬНЫЕ ВОПРОСЫ
% ============================================================
\chapter{Ответы на контрольные вопросы}

\section{Вопрос 1. Лемматизация и стемминг: сравнение}

\textit{Чем лемматизация отличается от стемминга? В каких сценариях предпочтительнее
использование каждого из методов применительно к юридическим текстам?}

\placeholder[Вставьте развёрнутый ответ (3–5 предложений)]{Ответ на вопрос 1}

\section{Вопрос 2. Стоп-слова и качество классификации}

\textit{Почему удаление стоп-слов улучшает качество классификации текстов? Можно
ли привести пример, когда удаление стоп-слов было бы вредным?}

\placeholder[Вставьте развёрнутый ответ (3–5 предложений)]{Ответ на вопрос 2}

\section{Вопрос 3. Ограничения ключевое-словарного классификатора}

\textit{Каковы принципиальные ограничения ключевое-словарного классификатора,
реализованного в Части 6? Какими методами можно преодолеть эти ограничения?}

\placeholder[Вставьте развёрнутый ответ (3–5 предложений)]{Ответ на вопрос 3}

\section{Вопрос 4. Специализированные библиотеки для русскоязычных текстов}

\textit{Почему для русскоязычных юридических текстов предпочтительно использование
специализированных библиотек, таких как \texttt{natasha}, вместо универсальных
решений?}

\placeholder[Вставьте развёрнутый ответ (3–5 предложений)]{Ответ на вопрос 4}

% ============================================================
%  РАЗДЕЛ 5. ИНДИВИДУАЛЬНОЕ ЗАДАНИЕ
% ============================================================
\chapter{Индивидуальное задание}

\section{Добавленные документы}

В данном разделе представлены результаты самостоятельного расширения корпуса
(не менее пяти новых документов).

\placeholder[Вставьте таблицу: ID\;|\;Категория\;|\;Источник\;|\;Объём (символов)]{Таблица добавленных документов}

\section{Обновлённые результаты классификации}

\placeholder[Вставьте обновлённую таблицу классификации всех документов]{Результаты классификации расширенного корпуса}

Обновлённая точность классификатора: \underline{\hspace{1.5cm}} /
\underline{\hspace{1cm}} = \underline{\hspace{1.5cm}}\%.

\section{Аналитический вывод}

\placeholder[Вставьте аналитический вывод объёмом не менее 150 слов: анализ изменений точности, выявленные закономерности, предложения по улучшению словаря]{Аналитический вывод по индивидуальному заданию}

% ============================================================
%  РАЗДЕЛ 6. ЗАКЛЮЧЕНИЕ
% ============================================================
\chapter{Заключение}

По итогам выполнения лабораторной работы были решены все поставленные задачи.
Реализован полный конвейер предобработки текстовых данных применительно к корпусу
судебных актов, включающий очистку текста, токенизацию, фильтрацию стоп-слов,
лемматизацию и извлечение именованных сущностей. На основе разработанного конвейера
построен ключевое-словарный классификатор и проведена его апробация на учебном корпусе.

\placeholder[Дополните данный абзац собственными выводами: что удалось, что вызвало затруднения, какие методы улучшили бы результаты]{Собственные выводы студента}

% ============================================================
%  СПИСОК ЛИТЕРАТУРЫ
% ============================================================
\chapter*{Список литературы}
\addcontentsline{toc}{chapter}{Список литературы}

\begin{thebibliography}{9}
  \bibitem{lec1}
    Материалы Лекции 1 по дисциплине «Информационно-поисковые системы».
    РУДН, 2025.
  \bibitem{lec2}
    Материалы Лекции 2 по дисциплине «Информационно-поисковые системы».
    РУДН, 2025.
  \bibitem{pymorphy2}
    Документация библиотеки \texttt{pymorphy2}.
    URL: \url{https://pymorphy2.readthedocs.io}
  \bibitem{natasha}
    Документация библиотеки \texttt{natasha}.
    URL: \url{https://natasha.github.io}
  \bibitem{nltk}
    Документация библиотеки NLTK (Natural Language Toolkit).
    URL: \url{https://www.nltk.org}
  \bibitem{kad}
    Электронная картотека арбитражных дел.
    URL: \url{https://kad.arbitr.ru}
\end{thebibliography}

% ============================================================
%  ПРИЛОЖЕНИЕ: СВОДНАЯ ВЕДОМОСТЬ БАЛЛОВ
% ============================================================
\chapter*{Приложение. Сводная ведомость баллов}
\addcontentsline{toc}{chapter}{Приложение. Сводная ведомость баллов}
\thispagestyle{plain}

\begin{tcolorbox}[infobox, title={Заполняется преподавателем}]
  \rowcolors{2}{LightBlue!30}{white}
  \begin{tabularx}{\linewidth}{X c c}
    \toprule
    \rowcolor{LightBlue}
    \textsf{\textbf{Критерий оценивания}} &
    \textsf{\textbf{Максимум}} &
    \textsf{\textbf{Получено}} \\
    \midrule
    Часть 2. Очистка текста (Задания 2.1–2.2)          & 20 & \\
    Часть 3. Токенизация и стоп-слова (Задания 3.1–3.2) & 15 & \\
    Часть 4. Лемматизация (Задания 4.1–4.2)             & 15 & \\
    Часть 5. Извлечение сущностей (Задания 5.1–5.2)     & 20 & \\
    Часть 6. Классификация (Задания 6.1–6.2)            & 15 & \\
    Задание 7.3 (повышенная сложность)                  & 15 & \\
    Индивидуальное задание                              & 15 & \\
    \midrule
    \textbf{ИТОГО}                                      & \textbf{100} & \\
    \bottomrule
  \end{tabularx}
\end{tcolorbox}

\vspace{24pt}
\noindent Подпись преподавателя: \underline{\hspace{6cm}}
\\[12pt]
\noindent Дата проверки: \underline{\hspace{4cm}}

\end{document}
